
\documentclass{article}
\usepackage{colortbl}
\usepackage{makecell}
\usepackage{multirow}
\usepackage{supertabular}

\begin{document}

\newcounter{utterance}

\centering \large Interaction Transcript for game `wordle', experiment `medium\_frequency\_words\_no\_clue\_no\_critic', episode 10 with glm{-}5.2.
\vspace{24pt}

{ \footnotesize  \setcounter{utterance}{1}
\setlength{\tabcolsep}{0pt}
\begin{supertabular}{c@{$\;$}|p{.15\linewidth}@{}p{.15\linewidth}p{.15\linewidth}p{.15\linewidth}p{.15\linewidth}p{.15\linewidth}}
    \# & \multicolumn{2}{c}{Player} && \multicolumn{2}{c}{Game Master} \\
    \hline

    \theutterance \stepcounter{utterance}  
    & & & \multicolumn{4}{p{0.6\linewidth}}{
        \cellcolor[rgb]{0.9,0.9,0.9}{
            \makecell[{{p{\linewidth}}}]{
                \texttt{\tiny{[P1$\langle$GM]}}
                \texttt{Ви — мовний чарівник, який любить вгадувати слова, використовуючи задані правила.} \\
\\ 
\texttt{Ласкаво просимо до Wordle! У вас є шість спроб, щоб вгадати цільове слово — допустиме слово з п'яти малих літер. Будь ласка, використовуйте теги "пояснення:" та "спроба:", щоб надати стислий опис для кожної спроби.} \\
\\ 
\texttt{Наприклад, якщо ваша спроба — "книга", ваша відповідь має бути} \\
\texttt{пояснення: це поширене слово з п'яти літер, і я починаю зі спроби з цього слова.} \\
\texttt{спроба: книга} \\
\\ 
\texttt{Після кожної спроби вашу відповідь буде перевірено, і ви отримаєте відгук, який вказує, які літери правильні (зелений), які літери правильні, але стоять не на своєму місці (жовтий), і які літери неправильні (червоний). Цей відгук може бути корисним для визначення, які літери слід включити або виключити у вашій наступній спробі.} \\
\\ 
\texttt{Наприклад, відгук для "книга" може бути:} \\
\texttt{відгук\_на\_спробу: к<жовтий> н<жовтий> и<зелений> г<жовтий> а<червоний>} \\
\\ 
\texttt{Пояснення має містити деталі про те, як відгук\_на\_спробу використовується, щоб дійти до нової спроби.} \\
\\ 
\texttt{Не генеруйте жодного іншого тексту.} \\
\\ 
\texttt{Почнімо з вашої першої спроби.} \\
            }
        }
    }
     \\ \\

    \theutterance \stepcounter{utterance}  
    & \multicolumn{4}{p{0.6\linewidth}}{
        \cellcolor[rgb]{0.561,0.737,0.561}{
            \makecell[{{p{\linewidth}}}]{
                \texttt{\tiny{[P1$\rangle$GM]}}
                \texttt{пояснення: Як мовний чарівник, я починаю з поширеного слова, що містить різноманітні літери, щоб швидко перевірити наявність голосних та приголосних.} \\
\texttt{спроба: слово} \\
            }
        }
    }
    & & \\ \\

    \theutterance \stepcounter{utterance}  
    & & & \multicolumn{2}{p{0.3\linewidth}}{
        \cellcolor[rgb]{0.9,0.9,0.9}{
            \makecell[{{p{\linewidth}}}]{
                \texttt{\tiny{[GM$|$GM]}}
                \texttt{attempts = 1} \\
\texttt{target\_word = балка} \\
\texttt{guess = слово} \\
\texttt{guess\_feedback = с<red> л<yellow> о<red> в<red> о<red>} \\
            }
        }
    }
    & & \\ \\

    \theutterance \stepcounter{utterance}  
    & & & \multicolumn{4}{p{0.6\linewidth}}{
        \cellcolor[rgb]{0.9,0.9,0.9}{
            \makecell[{{p{\linewidth}}}]{
                \texttt{\tiny{[P1$\langle$GM]}}
                \texttt{відгук\_на\_спробу: с<червоний> л<жовтий> о<червоний> в<червоний> о<червоний>} \\
\\ 
\texttt{Надайте відповідь лише в цьому форматі.} \\
\texttt{пояснення: деталі} \\
\texttt{спроба: слово} \\
            }
        }
    }
     \\ \\

    \theutterance \stepcounter{utterance}  
    & \multicolumn{4}{p{0.6\linewidth}}{
        \cellcolor[rgb]{0.561,0.737,0.561}{
            \makecell[{{p{\linewidth}}}]{
                \texttt{\tiny{[P1$\rangle$GM]}}
                \texttt{пояснення: Літера "л" є у слові, але не на другій позиції. Літери "с", "о", "в" відсутні. Шукаю нове слово без цих літер, але з "л" на іншій позиції.} \\
\texttt{спроба: кліпш} \\
            }
        }
    }
    & & \\ \\

    \theutterance \stepcounter{utterance}  
    & & & \multicolumn{2}{p{0.3\linewidth}}{
        \cellcolor[rgb]{0.9,0.9,0.9}{
            \makecell[{{p{\linewidth}}}]{
                \texttt{\tiny{[GM$|$GM]}}
                \texttt{Error: The guessed word is not a valid word for this game.} \\
            }
        }
    }
    & & \\ \\

    \theutterance \stepcounter{utterance}  
    & & & \multicolumn{4}{p{0.6\linewidth}}{
        \cellcolor[rgb]{0.9,0.9,0.9}{
            \makecell[{{p{\linewidth}}}]{
                \texttt{\tiny{[P1$\langle$GM]}}
                \texttt{Ваша спроба не є допустимим словом для цієї гри. Будь ласка, спробуйте ще раз.} \\
\\ 
\texttt{Надайте відповідь лише в цьому форматі.} \\
\texttt{пояснення: деталі} \\
\texttt{спроба: слово} \\
            }
        }
    }
     \\ \\

    \theutterance \stepcounter{utterance}  
    & \multicolumn{4}{p{0.6\linewidth}}{
        \cellcolor[rgb]{0.561,0.737,0.561}{
            \makecell[{{p{\linewidth}}}]{
                \texttt{\tiny{[P1$\rangle$GM]}}
                \texttt{пояснення: Літера "л" є у слові, але не на другій позиції. Літери "с", "о", "в" відсутні. Шукаю допустиме слово з п'яти літер без цих літер, але з "л" на іншій позиції.} \\
\texttt{спроба: ланець} \\
            }
        }
    }
    & & \\ \\

    \theutterance \stepcounter{utterance}  
    & & & \multicolumn{2}{p{0.3\linewidth}}{
        \cellcolor[rgb]{0.9,0.9,0.9}{
            \makecell[{{p{\linewidth}}}]{
                \texttt{\tiny{[GM$|$GM]}}
                \texttt{Error: The length of the guessed word is not 5.} \\
            }
        }
    }
    & & \\ \\

    \theutterance \stepcounter{utterance}  
    & & & \multicolumn{2}{p{0.3\linewidth}}{
        \cellcolor[rgb]{0.9,0.9,0.9}{
            \makecell[{{p{\linewidth}}}]{
                \texttt{\tiny{[GM$|$GM]}}
                \texttt{game\_result = ABORT} \\
            }
        }
    }
    & & \\ \\

\end{supertabular}
}

\end{document}
